\section*{Appendix outline}

We organize the appendices as follows. We devote Appendix~\ref{section:extension_hs} to the formalization of the extension of the results to Hilbert spaces.  Section~\ref{section:auxiliary_lemmas} presents auxiliary lemmata that are exploited in the remaining proofs. These are mostly analytic or simple probabilistic results that can be skipped on a first pass. Section~\ref{section:auxiliary_propositions} contains more involved technical propositions that are later combined to yield the proofs of Theorem~\ref{theorem:convergence_main_term} and Theorem~\ref{theorem:lower_order_term}. The proofs of such propositions are deferred to Appendix~\ref{appendix:auxiliary_propositions_proofs}.
Appendix~\ref{appendix:main_proofs} displays the proofs of the theoretical results exhibited in the main body of the paper, and Appendix~\ref{section:alternative_approaches} exhibits potential alternative approaches to the proposed empirical Bernstein inequality, illustrating the empirical benefits of the latter. %Lastly, Appendix~\ref{section:comparison} presents a comprehensive comparison with previous concentration inequalities.

%In Appendix~\ref{section:analysis_widths_std}, we prove via a short technical argument that the widths of the confidence intervals for the standard deviation also achieve $1/\sqrt{n}$ rates.

Throughout, we denote the probability space on which the random variables are defined by $\lp \Omega, \mathcal{F}, P \rp$. Furthermore, we use the standard asymptotic big-oh notations. Given functions $f$ and $g$, we write $f(n) = \bigO(g(n))$ if there exist constants $C, n_0 > 0$ such that $|f(n)| \leq C |g(n)|$ for all $n \geq n_0$. We write $f(n) = \widetilde{\bigO}(g(n))$ if $f(n) = \bigO(g(n) \, \mathrm{polylog}(n))$, where $\mathrm{polylog}(n)$ denotes a polylogarithmic factor in $n$. Finally, we use $f(n) = \Omega(g(n))$ to denote that there exist constants $c, n_0 > 0$ such that $f(n) \geq c g(n)$ for all $n \geq n_0$.